% !TEX program = xelatex
\documentclass[11pt,a4paper]{article}

\usepackage{fontspec}
\usepackage{polyglossia}
\usepackage{amsmath,amssymb,amsfonts}
\usepackage{xcolor}
\usepackage{graphicx}
\usepackage{geometry}

\geometry{left=2cm,right=2cm,top=2cm,bottom=2cm}

\setdefaultlanguage{persian}
\setotherlanguage{english}
\setmainfont{XB Niloofar}

\newcommand{\en}[1]{\textenglish{#1}}
\newcommand{\separator}{\vspace{1em}\noindent\rule{\linewidth}{0.4pt}\vspace{1em}}
\newcommand{\thickseparator}{\vspace{1em}\noindent\rule{\linewidth}{1pt}\vspace{1em}}

\title{\textbf{پاسخ سوال ۵} \\[0.4cm] \small تحلیل احساسات با استفاده از \en{Word2Vec Embeddings} و رگرسیون لجستیک}
\author{محمد باغبان زاده \\ شماره دانشجویی: شناسهٔ دانشجو}
\date{}
\linespread{1.35}

\begin{document}

\maketitle
\thickseparator

\section{مفاهیم پایه: ترکیب \en{Embeddings} و رگرسیون لجستیک برای تحلیل احساسات}

\subsection{نمایش کلمات با \en{Word2Vec Embeddings} (بخش \en{۵-۵} کتاب)}
در بخش \textbf{\en{۵-۵}} کتاب، مدل \en{Word2Vec} معرفی شده است که برای هر کلمه، یک بردار با ابعاد کم و متراکم (\en{dense vector}) تولید می‌کند. در این سوال، چنین بردارهایی برای کلمات زیر به صورت از پیش آموزش‌دیده داده شده است:

\begin{align*}
\text{\en{amazing}} &\rightarrow [0.7,\; 0.3,\; -0.1,\; 0.2] \\
\text{\en{film}} &\rightarrow [0.2,\; 0.1,\; 0.3,\; 0.1] \\
\text{\en{but}} &\rightarrow [0.0,\; 0.0,\; 0.0,\; 0.0] \quad \text{(کلمات توقف، بردار صفر)} \\
\text{\en{slow}} &\rightarrow [-0.4,\; 0.1,\; 0.2,\; -0.3] \\
\text{\en{boring}} &\rightarrow [-0.6,\; 0.0,\; 0.1,\; 0.1] \\
\text{\en{plot}} &\rightarrow [0.1,\; 0.2,\; 0.1,\; 0.0]
\end{align*}

\subsection{نمایش یک سند با میانگین بردارهای کلمات (بخش \en{۵-۳} کتاب)}
در بخش \textbf{\en{۵-۳}} کتاب (\en{Simple count-based embeddings})، روش میانگین گرفتن از بردارهای کلمات برای نمایش یک سند توضیح داده شده است. طبق صورت سوال، بردار سند به صورت \textbf{میانگین حسابی} بردارهای کلمات آن سند (پس از حذف کلمات توقف) محاسبه می‌شود.

\subsection{رده‌بندی با رگرسیون لجستیک (بخش \en{۴-۳} کتاب)}
در بخش \textbf{\en{۴-۳}} کتاب (\en{Classification with Logistic Regression})، فرمول محاسبه احتمال تعلق یک نمونه به کلاس مثبت به صورت زیر معرفی شده است:

\[
P(\text{\en{positive}}|x) = \sigma(w \cdot x + b) = \frac{1}{1 + e^{-(w \cdot x + b)}}
\]

که در آن \(x\) بردار ویژگی‌های نمونه (در اینجا بردار سند)، \(w\) بردار وزن‌های مدل و \(b\) مقدار بایاس است.

\thickseparator

\section{حل گام‌به‌گام مسأله}

\subsection{داده‌های مسأله}
\begin{align*}
\text{:جمله مورد نظر} &\quad \text{\en{``amazing film but slow boring plot''}} \\
\text{:بردار وزن‌های مدل} &\quad w = [0.6,\; -0.2,\; 0.3,\; 0.1] \\
\text{:مقدار بایاس} &\quad b = -0.05
\end{align*}

\subsection{گام ۱: حذف کلمات توقف (\en{Stop Words Removal})}
طبق صورت سوال، کلمه \en{but} (که بردار آن صفر است) به عنوان کلمه توقف محسوب می‌شود و باید حذف گردد. بنابراین کلمات مؤثر در جمله عبارتند از:

\[
\text{\en{amazing}},\; \text{\en{film}},\; \text{\en{slow}},\; \text{\en{boring}},\; \text{\en{plot}}
\]

\subsection{گام ۲: محاسبه بردار سند با روش میانگین حسابی}
برای هر بعد \(i\) از بردار سند، میانگین مقادیر آن بعد را در بین تمام کلمات محاسبه می‌کنیم:

\[
v_{\text{doc}}[i] = \frac{1}{n} \sum_{j=1}^{n} v_{\text{word}_j}[i]
\]

که در اینجا \(n = 5\) (تعداد کلمات پس از حذف کلمات توقف).

\textbf{بعد اول:}
\[
\frac{0.7 + 0.2 + (-0.4) + (-0.6) + 0.1}{5} = \frac{0.7 + 0.2 - 0.4 - 0.6 + 0.1}{5} = \frac{0.0}{5} = 0.0
\]

\textbf{بعد دوم:}
\[
\frac{0.3 + 0.1 + 0.1 + 0.0 + 0.2}{5} = \frac{0.7}{5} = 0.14
\]

\textbf{بعد سوم:}
\[
\frac{-0.1 + 0.3 + 0.2 + 0.1 + 0.1}{5} = \frac{0.6}{5} = 0.12
\]

\textbf{بعد چهارم:}
\[
\frac{0.2 + 0.1 + (-0.3) + 0.1 + 0.0}{5} = \frac{0.2 + 0.1 - 0.3 + 0.1 + 0.0}{5} = \frac{0.1}{5} = 0.02
\]

بنابراین بردار سند برابر است با:

\[
x = v_{\text{doc}} = [0.0,\; 0.14,\; 0.12,\; 0.02]
\]

\subsection{گام ۳: محاسبه مقدار خطی \(z = w \cdot x + b\)}

\[
w \cdot x = (0.6 \times 0.0) + (-0.2 \times 0.14) + (0.3 \times 0.12) + (0.1 \times 0.02)
\]

\[
w \cdot x = 0.0 + (-0.028) + 0.036 + 0.002 = 0.010
\]

\[
z = w \cdot x + b = 0.010 + (-0.05) = -0.04
\]

\subsection{گام ۴: اعمال تابع سیگموئید برای محاسبه احتمال}

\[
P(\text{\en{positive}}|x) = \sigma(z) = \frac{1}{1 + e^{-z}} = \frac{1}{1 + e^{-(-0.04)}} = \frac{1}{1 + e^{0.04}}
\]

مقدار \(e^{0.04}\) را تقریب می‌زنیم (\(e^{0.04} \approx 1.0408\)):

\[
P(\text{\en{positive}}|x) = \frac{1}{1 + 1.0408} = \frac{1}{2.0408} \approx 0.4900
\]

\separator
\section{جمع‌بندی و پاسخ نهایی}

\begin{itemize}
    \item \textbf{بردار سند (پس از میانگین گرفتن):} \(x = [0.0,\; 0.14,\; 0.12,\; 0.02]\)
    \item \textbf{مقدار خطی \(z = w \cdot x + b\):} \(z = -0.04\)
    \item \textbf{احتمال مثبت بودن سند:} \(P(\text{\en{positive}}|x) \approx 0.4900\)
\end{itemize}

\thickseparator
\begin{center}
\fbox{\parbox{0.85\linewidth}{\centering
\textbf{پاسخ نهایی:} \\
\[
\boxed{P(\text{positive} | \text{amazing film but slow boring plot}) \approx 0.4900}
\]
}}
\end{center}

\thickseparator

\section{تفسیر نتیجه}
مقدار احتمال محاسبه‌شده \((\approx 0.49)\) کمی کمتر از آستانه \(0.5\) است. بنابراین اگر بخواهیم با آستانه \(0.5\) تصمیم‌گیری کنیم (طبق بخش \en{۴-۳} کتاب)، این سند به کلاس \textbf{منفی (\en{negative})} تعلق خواهد گرفت. این نتیجه با توجه به حضور کلمات منفی مانند \en{slow} و \en{boring} در جمله، منطقی به نظر می‌رسد.

\thickseparator

\end{document}